\documentclass{homework}

\usepackage{enumitem}
\usepackage{tcolorbox}

\usepackage{tikz}
\usetikzlibrary{positioning,calc}

\usepackage{listings}

\lstset{
basicstyle=\small\ttfamily,
columns=flexible,
breaklines=true
}

\usepackage{fancyvrb}
\fvset{listparameters=\setlength{\topsep}{0pt}\setlength{\partopsep}{0pt}}

\title{03}

\begin{document}
\maketitle

\textbf{Organisational details}

\begin{itemize}
  \item Include in your report approximately how long the assignment took you.
  \item You must follow the submission instructions outlined on GitLab.
  The theory part must be formatted as Markdown.
  \item The GitLab pipeline must succeed, or I will not grade your assignment (you will receive 0).
  \item I will provide feedback on GitLab and/or Discord.
  Check both after receiving your grade.
  \begin{itemize}
    \item My initial feedback may be quite minimal as there are so many of you.
    However, you can request additional feedback.
    Specific questions are also welcome.
    \item Do not share the feedback you receive, as this may give others an unfair advantage.
  \end{itemize}
  \item You may share hints with each other, but not answers or code.
  Ask questions in the Discord server if they could benefit everyone.
  \item You may -- and I encourage you to -- ask me questions.
  \item If I ask, you must be able to convince me that the homework is your own work.
  Otherwise, I will report you to the faculty and fail you.
\end{itemize}

\newpage

\begin{center}
  Theory tasks
\end{center}

\textbf{Additional instructions}

\begin{itemize}
  \item \textbf{Please do not use AI for the theoretical part.}
  Challenge your understanding and look for additional written sources if the lecture and provided materials are not enough.
  Critical thinking and clear communication are skills -- use them or lose them.
  \item You must reference all external materials you rely on, including AI tools, if you choose to use them anyway.
\end{itemize}

\begin{task}{ROCA}
  Read \href{https://crocs.fi.muni.cz/public/papers/rsa_ccs17}{\emph{this blog post}} by the researchers who discovered ROCA.
  More than 750 000 Estonian ID cards were impacted by this vulnerability.

  Answer the following:
  \begin{itemize}
    \item What caused the vulnerability?
    Was it a weak random number generator?
    \item Did the researchers break 2048-bit RSA in general?
    What was the cost of breaking a single vulnerable 2048-bit key?\footnote{Note that it would likely be more affordable today.}
    \item How can you check whether your key is affected?
    \item Estonia's response was to switch ID card keys from RSA to ECC.
    Why did this work?
    \item The researchers also suggest generating the key outside the device and importing it.
    Why is this a poor option for an ID card?
    (Hint: look up why the PPA sued Gemalto in 2018.)
  \end{itemize}
\end{task}

\begin{task}{Padding oracles and Bleichenbacher}
  Read the introduction of \href{https://en.wikipedia.org/wiki/Padding_oracle_attack}{\textit{Padding oracle attack}} and skim the rest of the article.
  Then read through \href{https://www.shysecurity.com/post/20160622-Bleichenbacher%20Attack}{\textit{Bleichenbacher Chosen Ciphertext Attack}}.
  Supporting materials (optional):
  \begin{itemize}
    \item \url{https://crypto.stackexchange.com/a/37480}
    \item \href{https://wiki.elvis.science/index.php?title=Bleichenbacher_Attack#Steps_of_the_Attack}{\textit{Bleichenbacher Attack}} by the ELVIS project -- steps of Bleichenbacher's attack
    \item \href{https://robertheaton.com/2013/07/29/padding-oracle-attack/}{\textit{The padding oracle attack}} by Robert Heaton -- description of the CBC padding oracle
  \end{itemize}

  Answer the following:
  \begin{itemize}
    \item What is a padding oracle attack in general, i.e. not specific to RSA?
    Does the attacker need the capabilities of an IND-CCA1 or an IND-CCA2 adversary?
    Why?
    \item Without going into the mathematical details, how do the oracle's answers enable plaintext recovery in Bleichenbacher's attack?
    \item RSAES-OAEP is proven IND-CCA2 secure\footnote{In the random oracle model (ROM).}.
    Does that mean that in practice, RSAES-OAEP implementations cannot be vulnerable to such attacks?
    Explain why or why not.
  \end{itemize}
\end{task}

\begin{task}{Opening up}
  Recall that for some secret $x\getsu\ZZ_q$ and a generator $G$, the ElGamal public key is $H\gets xG$.
  Encryption is then carried out with:
  \[
    r \getsu \ZZ_q, \ENC_H(M; r) = (rG, M + rH) = (U, V) = c
  \]
  where $r$ is the randomness which makes ElGamal a probabilistic encryption scheme.

  Suppose now that $r$ is leaked, meaning that you know the $r$ used during encryption.
  Are you then able to recover $M$ from $c$ without the secret key?
  If yes, how so?

  Can you also decrypt other messages (without gaining access to new values of $r$)?
\end{task}

\begin{task}{Following up}
  Regardless of the public-key encryption scheme, the encryption randomness must always be freshly sampled, kept secret, and discarded after encryption.

  Explain how, if the encryption randomness leaks and the set of possible messages is known (e.g. the choices \enquote{yes} or \enquote{no} in a referendum), an adversary can recover the message, regardless of the public-key encryption scheme used (e.g. RSAES-OAEP or lifted ElGamal).
\end{task}

\newpage
\setcounter{task}{0}

\begin{center}
  Practical tasks
\end{center}

\textbf{Additional instructions}

\begin{itemize}
  \item I must be able to run your program with Python 3.14 and OpenSSL 4.0.3\footnotemark{}.
  \footnotetext{You may use older versions of Python and OpenSSL as long as they are interoperable with mine.}
  \item You may not import any third-party module.
  You are free to import any built-in module.
  If you feel the need to import a third-party module, please check with me beforehand!
  \item All errors must be handled: print a descriptive message starting with \enquote{Handled:} and exit with error code \texttt{1}.
  \item You must identify potential edge cases yourself.
  \item You may use AI as assistance, e.g. to check your understanding or to ask how to loop in Python.
  You may not use AI to generate code logic, library calls, or cryptographic operations, nor to identify or implement edge cases.
  If you use AI in any form, state in your report where and what you used it for.
  \item You may change anything in the template files, unless specified otherwise, as long as the GitLab pipeline succeeds.
\end{itemize}

\begin{task}{RSA interop}
  Homework template file: \texttt{ossl\_rsa\_template.py}

  Create a Python 3 program that performs RSAES-OAEP encryption and decryption by calling OpenSSL.
  Recall that RSA needs padding for probabilistic encryption, and that OAEP in particular prevents attacks such as Bleichenbacher's.
  Use SHA-256 as the OAEP hash function.

  Your program must provide three separate tools:
  \begin{itemize}
    \item \texttt{genpkey} -- generate an RSA private key

    Your program must generate a $3072$-bit RSA private key with OpenSSL and export it into a file in PEM format.
    Your program \emph{must not} print the key-generation progress bar to the console.

    The private key must be password protected according to the following:
    \begin{itemize}
      \item The key derivation function for obtaining the encryption key must be PBKDF2, and the parameters must be:
      \begin{itemize}
        \item \texttt{HMAC-SHA512} for the PRF,
        \item an iteration count of $210000$ (OWASP recommendation).
      \end{itemize}
      \item The key must be encrypted with AES-256 in CBC mode and use the PKCS\#8 format.
      \item At no point should the private key be written unencrypted to the disk.
    \end{itemize}

    \begin{tcolorbox}[title=Note]
      Since version 3.0, OpenSSL uses PKCS\#8 by default when exporting an encrypted private key in PEM, but \texttt{genpkey} and \texttt{pkey} do not allow for the fine-grained control needed for this task.

      Hint: can you find a tool in the OpenSSL documentation which allows more fine-grained control over PKCS\#8 encrypted private keys?
    \end{tcolorbox}

    You can see what the expected private key structure looks like by dumping the ASN.1 structure of the example private key \texttt{rsa.demo.priv.pem}.

    Example invocation:
    \begin{Verbatim}
ossl_rsa.py genpkey --out key.pem --pass secret
    \end{Verbatim}

    \item \texttt{pkey} -- public key derivation
    
    Your program must extract the public key from a private key and export it into a file in PEM format by using OpenSSL.

    Example invocation:
    \begin{Verbatim}
ossl_rsa.py pkey --in key.pem \
    --out pub.pem --pass secret
    \end{Verbatim}

    \item \texttt{pkeyutl} -- file encryption and decryption
    
    Your program must output the encryption or decryption of the specified input file with the specified key to the specified output file using OpenSSL.

    Example invocation for encryption:
    \begin{Verbatim}
ossl_rsa.py pkeyutl --encrypt \
    --inkey pub.pem --in message.txt --out data.enc
    \end{Verbatim}

    Example invocation for decryption:
    \begin{Verbatim}
ossl_rsa.py pkeyutl \
    --inkey key.pem --in data.enc --out data.dec \
    --pass secret
    \end{Verbatim}
  \end{itemize}

  Your program must not crash if parameters are missing, the wrong password is provided, OpenSSL fails, etc.
  Some error handling and hints are part of the template, but not all of it.

  You do not have to (but you may) handle file I/O errors, i.e. you may assume that input files exist and output files can be written to.

  \begin{tcolorbox}[title=Reminder]
    In practice, providing secrets (e.g. keys, passwords, etc.) directly on the CLI as an argument to a program is a terrible idea, since they typically get logged into your terminal history!
  \end{tcolorbox}
\end{task}

\end{document}
